% !TeX root = ../main.tex

\begin{cnabstract}
本文以申万宏源研究中的柏顿·墨基尔选股法为复现对象，保留基础复现与同规则延展，并建立四条有明确对照关系的优化路径：正PE与初始五年修正、全截面独立排名，以及分别叠加原有仓位管理的两条路径。原始时期的补充检验中，151个月敏感性口径下P1、P2年化分别为38.77\%、40.97\%，高于同窗口基础版36.44\%；两条风控路径以降低收益换取回撤收窄。基础复现实际覆盖2003年4月至2015年6月147个月；完整延展及四路径统一覆盖至2026年9月，共282个月。结果显示，基础延展年化收益为26.17\%，最大回撤为$-$70.11\%；路径一、二年化收益分别为29.64\%和27.38\%，最大回撤分别为$-$68.87\%和$-$64.15\%。路径三、四将最大回撤收窄至$-$32.74\%和$-$34.86\%，但年化收益降至23.19\%和20.85\%。路径一在全期收益上居前，路径三在所设零利率夏普与回撤指标上较优，但分段排序不稳定；2025—2026年追加区间中，基础延展优于四条路径。机制核查表明，排名母体会改变补位池及最终核心保留，风控路径未出现零仓位月份，不能将风险改善归因于已经生效的止损清仓。四路径仍沿用报告期财务映射、月度再等权和费用近似，独立CSMAR重建逐股复现基准，282个月最大收益差约$5\times10^{-16}$；与原报告的差异可定位到比较区间、调仓安排及部分历史缺口。本文据此将其解释为具有可复核计算依据的历史规则比较，而非已证明的可交易超额收益。
\end{cnabstract}
\cnkeywords{墨基尔选股法；持续盈余成长；低市盈率；排名母体；波动率目标；策略复现}
